% ============================================================
\chapter{Estimation Bay\'esienne et Fonctions de Perte}
\label{ch:estimation}
% ============================================================

En statistique fr\'equentiste, un estimateur est \og bon \fg s'il est sans biais et de variance minimale. En statistique bay\'esienne, la question est formul\'ee diff\'eremment~: quel r\'esum\'e de la loi a posteriori choisir, et \`a quel co\^ut~? La r\'eponse d\'epend de ce que l'on consid\`ere comme une erreur grave. Se tromper de beaucoup est-il proportionnellement pire que se tromper de peu (perte quadratique), ou toute erreur au-del\`a d'un seuil est-elle \'egalement inacceptable (perte 0-1)~? Abraham Wald, dans ses travaux fondateurs des ann\'ees 1940 sur la th\'eorie de la d\'ecision, a montr\'e que cette question de la \emph{fonction de perte} est in\'evitable~: tout estimateur est, consciemment ou non, la solution d'un probl\`eme de d\'ecision.

Ce chapitre explore ce lien entre estimation et d\'ecision, et montre comment chaque fonction de perte conduit \`a un estimateur optimal diff\'erent~: la moyenne a posteriori, la m\'ediane a posteriori, le mode a posteriori.

\begin{intuition}[title=\textbf{Id\'ee centrale}]
L'estimation bay\'esienne ne se r\'eduit pas \`a un estimateur ponctuel
unique~: elle fournit toute la loi a posteriori. N\'eanmoins, lorsqu'un
r\'esum\'e ponctuel est n\'ecessaire, le choix de l'estimateur optimal
d\'epend de la \textbf{fonction de perte} retenue.
\end{intuition}

% ------------------------------------------------------------
\section{Th\'eorie de la d\'ecision bay\'esienne}
% ------------------------------------------------------------

\begin{definition}[Probl\`eme de d\'ecision]
Un \textbf{probl\`eme de d\'ecision} est un triplet $(\Theta, \mathcal{A}, L)$
o\`u~:
\begin{itemize}
  \item $\Theta$ est l'espace des param\`etres,
  \item $\mathcal{A}$ est l'espace des actions (d\'ecisions),
  \item $L : \Theta \times \mathcal{A} \to \R^+$ est la fonction de perte.
\end{itemize}
\end{definition}

\begin{definition}[Risque a posteriori]
Le \textbf{risque a posteriori} (ou perte esp\'er\'ee a posteriori)
d'une action $a \in \mathcal{A}$ est~:
\[
  \rho(\pi, a \mid x) = \E_{\theta \mid x}[L(\theta, a)]
  = \int_\Theta L(\theta, a)\,\pi(\theta \mid x)\,d\theta.
\]
\end{definition}

\begin{definition}[Estimateur de Bayes]
L'\textbf{estimateur de Bayes} sous la perte $L$ est l'action qui
minimise le risque a posteriori~:
\[
  \hat{\theta}_B(x) = \argmin_{a \in \mathcal{A}} \rho(\pi, a \mid x).
\]
\end{definition}

\begin{definition}[Risque bay\'esien (int\'egr\'e)]
Le \textbf{risque bay\'esien} d'un estimateur $\delta(x)$ est~:
\[
  r(\pi, \delta) = \E_x\left[\rho(\pi, \delta(x) \mid x)\right]
  = \int \int L(\theta, \delta(x))\,f(x \mid \theta)\,\pi(\theta)\,dx\,d\theta.
\]
\end{definition}

\begin{theoreme}[Optimalit\'e de l'estimateur de Bayes]
L'estimateur de Bayes $\hat{\theta}_B$ minimise le risque bay\'esien
$r(\pi, \delta)$ parmi tous les estimateurs.
\end{theoreme}

\begin{proof}
Le risque bay\'esien peut s'\'ecrire~:
\[
  r(\pi, \delta) = \int \rho(\pi, \delta(x) \mid x)\,m(x)\,dx,
\]
o\`u $m(x)$ est la vraisemblance marginale. Comme $m(x) \geq 0$,
minimiser $r(\pi, \delta)$ revient \`a minimiser $\rho(\pi, \delta(x) \mid x)$
pour chaque $x$, ce qui donne $\hat{\theta}_B(x)$.
\end{proof}

% ------------------------------------------------------------
\section{Fonctions de perte classiques}
% ------------------------------------------------------------

\begin{formules}[title=\textbf{Estimateurs de Bayes selon la perte}]
\begin{center}
\renewcommand{\arraystretch}{1.5}
\begin{tabular}{p{4.5cm} p{4cm} p{5cm}}
\toprule
\textbf{Fonction de perte} & \textbf{Expression} & \textbf{Estimateur de Bayes} \\
\midrule
Quadratique & $L(\theta, a) = (\theta - a)^2$ &
  $\E[\theta \mid x]$ (moyenne post.) \\
Valeur absolue & $L(\theta, a) = \abs{\theta - a}$ &
  $\text{Med}[\theta \mid x]$ (m\'ediane post.) \\
0--1 (tout ou rien) & $L(\theta, a) = \mathbf{1}(\abs{\theta - a} > \varepsilon)$ &
  Mode$[\theta \mid x]$ (MAP) \\
LINEX & $L(\theta, a) = e^{c(\theta-a)} - c(\theta-a) - 1$ &
  $\frac{1}{c}\log \E[e^{c\theta} \mid x]$ \\
Quadratique pond\'er\'ee & $L(\theta, a) = w(\theta)(\theta - a)^2$ &
  $\frac{\E[w(\theta)\theta \mid x]}{\E[w(\theta) \mid x]}$ \\
\bottomrule
\end{tabular}
\end{center}
\end{formules}

\subsection{Perte quadratique}

\begin{theoreme}[Estimateur de Bayes sous perte quadratique]
\label{thm:bayes_quad}
Sous la perte quadratique $L(\theta, a) = (\theta - a)^2$,
l'estimateur de Bayes est la moyenne a posteriori~:
\[
  \hat{\theta}_B(x) = \E[\theta \mid x].
\]
\end{theoreme}

\begin{proof}
\begin{align*}
  \frac{\partial}{\partial a} \E[(\theta - a)^2 \mid x]
  &= -2\,\E[\theta - a \mid x]
  = -2\big(\E[\theta \mid x] - a\big).
\end{align*}
Ce qui s'annule pour $a = \E[\theta \mid x]$. La d\'eriv\'ee seconde
est $2 > 0$, confirmant un minimum.
\end{proof}

\begin{corollaire}
Le risque a posteriori sous perte quadratique est la variance a
posteriori~:
\[
  \rho(\pi, \hat{\theta}_B \mid x) = \Var[\theta \mid x].
\]
\end{corollaire}

\subsection{Perte valeur absolue}

\begin{theoreme}[Estimateur de Bayes sous perte absolue]
Sous la perte $L(\theta, a) = \abs{\theta - a}$, l'estimateur de Bayes
est la m\'ediane a posteriori.
\end{theoreme}

\begin{proof}
La d\'eriv\'ee g\'en\'eralis\'ee du risque a posteriori est~:
\[
  \frac{\partial}{\partial a} \E[\abs{\theta - a} \mid x]
  = \PP(\theta < a \mid x) - \PP(\theta > a \mid x),
\]
qui s'annule quand $\PP(\theta < a \mid x) = 1/2$, i.e., $a$ est la
m\'ediane de $\pi(\theta \mid x)$.
\end{proof}

\subsection{Maximum a posteriori (MAP)}

\begin{definition}[Estimateur MAP]
L'\textbf{estimateur du maximum a posteriori} (MAP) est~:
\[
  \hat{\theta}_{\text{MAP}} = \argmax_\theta \pi(\theta \mid x)
  = \argmax_\theta \big[\log L(\theta \mid x) + \log \pi(\theta)\big].
\]
\end{definition}

\begin{remarque}
Le MAP co\"incide avec le MLE quand le prior est uniforme. Avec un
prior gaussien $\mathcal{N}(0, \tau^2)$, le MAP correspond \`a la
r\'egularisation de Ridge~:
\[
  \hat{\theta}_{\text{MAP}} = \argmin_\theta \left[
  -\log L(\theta \mid x) + \frac{\norm{\theta}^2}{2\tau^2}\right].
\]
\end{remarque}

\begin{proposition}[MAP et r\'egularisation Lasso]
Si le prior est un Laplace $\pi(\theta) \propto e^{-\lambda \abs{\theta}}$,
le MAP correspond \`a la r\'egularisation Lasso ($\ell_1$).
\end{proposition}

% ------------------------------------------------------------
\section{Intervalles de cr\'edibilit\'e}
% ------------------------------------------------------------

\begin{definition}[Intervalle de cr\'edibilit\'e]
Un \textbf{intervalle de cr\'edibilit\'e} \`a $100(1-\alpha)\%$ est un
intervalle $[a, b]$ tel que~:
\[
  \PP(\theta \in [a, b] \mid x) = 1 - \alpha.
\]
\end{definition}

\begin{definition}[HDI --- Highest Density Interval]
L'\textbf{intervalle de plus haute densit\'e} (HDI) est l'intervalle
de cr\'edibilit\'e de longueur minimale~:
\[
  \text{HDI}_{1-\alpha} = \{
  \theta : \pi(\theta \mid x) \geq c_\alpha\},
\]
o\`u $c_\alpha$ est choisi tel que
$\PP(\pi(\theta \mid x) \geq c_\alpha \mid x) = 1 - \alpha$.
\end{definition}

\begin{attention}[title=\textbf{Cr\'edibilit\'e vs. confiance}]
Un intervalle de cr\'edibilit\'e \`a $95\%$ signifie que $\theta$
appartient \`a l'intervalle avec probabilit\'e $0.95$ \textit{conditionnellement
aux donn\'ees observ\'ees}. Un intervalle de confiance \`a $95\%$
signifie que la proc\'edure contient $\theta$ dans $95\%$ des
r\'ep\'etitions \textit{hypoth\'etiques}.
\end{attention}

\begin{exemple}
Pour le mod\`ele Bernoulli--B\^eta avec $n = 100$, $s = 30$,
$\theta \sim \text{Be}(1,1)$, le posterior est $\text{Be}(31, 71)$.

L'intervalle de cr\'edibilit\'e \'equi-queue \`a $95\%$ est~:
\[
  [F^{-1}_{31,71}(0.025),\; F^{-1}_{31,71}(0.975)]
  \approx [0.216,\; 0.397].
\]
\end{exemple}

% ------------------------------------------------------------
\section{Comparaison bay\'esien--fr\'equentiste}
% ------------------------------------------------------------

\begin{definition}[Risque fr\'equentiste]
Le risque fr\'equentiste d'un estimateur $\delta(x)$ est~:
\[
  R(\theta, \delta) = \E_{x \mid \theta}[L(\theta, \delta(x))].
\]
\end{definition}

\begin{definition}[Admissibilit\'e]
Un estimateur $\delta$ est \textbf{admissible} s'il n'existe pas
d'estimateur $\delta'$ tel que $R(\theta, \delta') \leq R(\theta, \delta)$
pour tout $\theta$, avec in\'egalit\'e stricte pour au moins un $\theta$.
\end{definition}

\begin{theoreme}[Admissibilit\'e des estimateurs de Bayes]
Sous des conditions de r\'egularit\'e, tout estimateur de Bayes dont le
risque bay\'esien est fini est admissible.
\end{theoreme}

\begin{theoreme}[Classe compl\`ete]
Sous des conditions g\'en\'erales, la classe des estimateurs de Bayes
(propres et limites) forme une \textbf{classe compl\`ete}~: tout
estimateur admissible est un estimateur de Bayes (\'eventuellement
g\'en\'eralis\'e).
\end{theoreme}

\begin{exemple}
\textbf{Estimateur de James--Stein.} Pour $X \sim \mathcal{N}_p(\theta, I_p)$
avec $p \geq 3$, l'estimateur $\hat{\theta}_{\text{JS}} = (1 - \frac{p-2}{\norm{X}^2})X$
domine le MLE $\hat{\theta} = X$ sous perte quadratique. C'est un
estimateur de Bayes g\'en\'eralis\'e.
\end{exemple}

% ------------------------------------------------------------
\section{Estimation ponctuelle~: exemples}
% ------------------------------------------------------------

\begin{exemple}
\textbf{Mod\`ele Poisson--Gamma.}
Avec $X_i \sim \text{Poi}(\lambda)$ et $\lambda \sim \text{Ga}(\alpha, \beta)$,
le posterior est $\text{Ga}(\alpha + \sum x_i, \beta + n)$. Les
estimateurs ponctuels sont~:
\begin{align*}
  \hat{\lambda}_{\text{moyenne}} &= \frac{\alpha + \sum x_i}{\beta + n}, \\
  \hat{\lambda}_{\text{MAP}} &= \frac{\alpha + \sum x_i - 1}{\beta + n}
  \quad (\text{si } \alpha + \sum x_i > 1), \\
  \hat{\lambda}_{\text{m\'ediane}} &\approx \hat{\lambda}_{\text{moyenne}}
  - \frac{1}{3(\beta + n)} \quad (\text{approx. pour grande } \alpha').
\end{align*}
\end{exemple}

\begin{exemple}
\textbf{Mod\`ele Normal--Normal.}
Avec $X_i \sim \mathcal{N}(\mu, \sigma^2)$, $\sigma^2$ connu, et
$\mu \sim \mathcal{N}(\mu_0, \tau_0^2)$~: le posterior est gaussien
donc sym\'etrique, et les trois estimateurs co\"incident~:
\[
  \hat{\mu}_{\text{moyenne}} = \hat{\mu}_{\text{MAP}} = \hat{\mu}_{\text{m\'ediane}}
  = \mu_n = \tau_n^2\left(\frac{\mu_0}{\tau_0^2} + \frac{n\bar{x}}{\sigma^2}\right).
\]
\end{exemple}

% ------------------------------------------------------------
\section{Perte LINEX et estimation asym\'etrique}
% ------------------------------------------------------------

\begin{definition}[Perte LINEX]
La perte LINEX (\textit{LINear-EXponential}) est~:
\[
  L(\theta, a) = e^{c(\theta - a)} - c(\theta - a) - 1, \qquad c \neq 0.
\]
Pour $c > 0$, la surestimation est p\'enalis\'ee exponentiellement.
Pour $c < 0$, c'est la sous-estimation.
\end{definition}

\begin{theoreme}[Estimateur de Bayes sous perte LINEX]
L'estimateur de Bayes sous perte LINEX est~:
\[
  \hat{\theta}_{\text{LINEX}} = \frac{1}{c}\log \E[e^{c\theta} \mid x].
\]
\end{theoreme}

\begin{proof}
Le risque a posteriori est~:
\[
  \rho(a) = \E[e^{c(\theta-a)} \mid x] - c\,\E[\theta - a \mid x] - 1.
\]
En d\'erivant par rapport \`a $a$ et en annulant~:
\[
  -c\,e^{-ca}\,\E[e^{c\theta} \mid x] + c = 0
  \implies e^{ca} = \E[e^{c\theta} \mid x],
\]
d'o\`u $a = \frac{1}{c}\log \E[e^{c\theta} \mid x]$.
\end{proof}

% ------------------------------------------------------------
\section{R\'etr\'ecissement bay\'esien (\textit{shrinkage})}
% ------------------------------------------------------------

\begin{proposition}[R\'etr\'ecissement vers la moyenne a priori]
Dans le mod\`ele Normal--Normal, la moyenne a posteriori s'\'ecrit~:
\[
  \E[\mu \mid x] = (1 - B)\bar{x} + B\mu_0,
  \qquad B = \frac{\sigma^2/n}{\tau_0^2 + \sigma^2/n},
\]
o\`u $B \in (0,1)$ est le \textbf{facteur de r\'etr\'ecissement}.
L'estimateur bay\'esien est r\'etr\'eci vers la moyenne a priori $\mu_0$.
\end{proposition}

% ------------------------------------------------------------
\section{Impl\'ementation Python}
% ------------------------------------------------------------

\begin{codeblock}[title=\textbf{Estimation bay\'esienne et intervalles de cr\'edibilit\'e}]
\begin{minted}{python}
import numpy as np
from scipy import stats
import matplotlib.pyplot as plt

# Données Bernoulli
np.random.seed(42)
n, s = 100, 30

# Posterior Beta(31, 71)
a_post, b_post = 1 + s, 1 + n - s
posterior = stats.beta(a_post, b_post)

# Estimateurs ponctuels
mean_post = posterior.mean()
median_post = posterior.median()
mode_post = (a_post - 1) / (a_post + b_post - 2)

print(f"Moyenne a posteriori: {mean_post:.4f}")
print(f"Médiane a posteriori: {median_post:.4f}")
print(f"MAP: {mode_post:.4f}")
print(f"MLE: {s/n:.4f}")

# Intervalles de crédibilité
ci_equi = posterior.ppf([0.025, 0.975])
print(f"IC équi-queue 95%: [{ci_equi[0]:.4f}, {ci_equi[1]:.4f}]")

# HDI (via ArviZ)
import arviz as az
samples = posterior.rvs(100000)
hdi = az.hdi(samples, hdi_prob=0.95)
print(f"HDI 95%: [{hdi[0]:.4f}, {hdi[1]:.4f}]")
\end{minted}
\end{codeblock}

\begin{codeblock}[title=\textbf{Comparaison des fonctions de perte}]
\begin{minted}{python}
# Visualisation des fonctions de perte
delta = np.linspace(-3, 3, 500)

fig, axes = plt.subplots(1, 3, figsize=(14, 4))

# Perte quadratique
axes[0].plot(delta, delta**2)
axes[0].set_title("Perte quadratique")
axes[0].set_xlabel(r"$\theta - a$")

# Perte valeur absolue
axes[1].plot(delta, np.abs(delta))
axes[1].set_title("Perte valeur absolue")
axes[1].set_xlabel(r"$\theta - a$")

# Perte LINEX (c=1 et c=-1)
for c in [1, -1]:
    loss = np.exp(c*delta) - c*delta - 1
    axes[2].plot(delta, loss, label=f"c={c}")
axes[2].set_title("Perte LINEX")
axes[2].set_xlabel(r"$\theta - a$")
axes[2].legend()

for ax in axes:
    ax.set_ylabel("Perte")

plt.tight_layout()
plt.savefig("loss_functions.pdf")
\end{minted}
\end{codeblock}

\begin{codeblock}[title=\textbf{Effet de r\'etr\'ecissement bay\'esien}]
\begin{minted}{python}
import pymc as pm
import arviz as az

# Simulation: J groupes, estimation simultanée
np.random.seed(0)
J = 8
true_mu = np.array([-2, -1, 0, 0.5, 1, 1.5, 2, 3])
sigma = 1.0
n_per_group = 5

data = [np.random.normal(mu, sigma, n_per_group)
        for mu in true_mu]
x_bar = np.array([d.mean() for d in data])

# Estimateurs MLE vs Bayes (prior N(0, 4))
tau0 = 2.0
B = (sigma**2 / n_per_group) / (tau0**2 + sigma**2 / n_per_group)
bayes_est = (1 - B) * x_bar + B * 0  # mu_0 = 0

fig, ax = plt.subplots(figsize=(8, 6))
ax.scatter(range(J), true_mu, marker="*", s=150,
           label="Vrai $\\mu_j$", zorder=3)
ax.scatter(range(J), x_bar, marker="o",
           label="MLE $\\bar{x}_j$")
ax.scatter(range(J), bayes_est, marker="s",
           label="Bayes (shrinkage)")
ax.axhline(0, color="gray", linestyle=":", alpha=0.5)
ax.set_xlabel("Groupe $j$")
ax.set_ylabel(r"$\mu_j$")
ax.legend()
ax.set_title("Rétrécissement bayésien")
plt.tight_layout()
plt.savefig("shrinkage.pdf")
\end{minted}
\end{codeblock}

% ------------------------------------------------------------
\section{Exercices}
% ------------------------------------------------------------

\begin{exercice}
D\'emontrer que l'estimateur de Bayes sous la perte pond\'er\'ee
$L(\theta, a) = w(\theta)(\theta - a)^2$ est~:
\[
  \hat{\theta}_B = \frac{\E[w(\theta)\theta \mid x]}{\E[w(\theta) \mid x]}.
\]
\end{exercice}

\begin{exercice}
Pour le mod\`ele Poisson--Gamma, calculer l'estimateur de Bayes sous
la perte LINEX avec $c = 1$ et comparer avec la moyenne a posteriori.
\end{exercice}

\begin{exercice}
Montrer que dans le cas gaussien unidimensionnel, le HDI co\"incide
avec l'intervalle \'equi-queue.
\end{exercice}

\begin{exercice}
Soit $X \sim \mathcal{N}_p(\theta, I_p)$ avec $\theta \sim \mathcal{N}_p(0, \tau^2 I_p)$.
\begin{enumerate}
  \item Calculer l'estimateur de Bayes sous perte quadratique.
  \item Montrer qu'il s'agit d'un estimateur de r\'etr\'ecissement
        de la forme $c \cdot X$ avec $c \in (0,1)$.
  \item Comparer avec l'estimateur de James--Stein.
\end{enumerate}
\end{exercice}

\begin{exercice}
\textbf{(Num\'erique)} Simuler $n = 50$ observations de
$\mathcal{N}(3, 4)$. Comparer les estimateurs bay\'esiens de $\mu$
(moyenne, m\'ediane, MAP) avec le MLE pour diff\'erents priors~:
$\mathcal{N}(0, 100)$ (vague), $\mathcal{N}(0, 1)$ (informatif
centr\'e), $\mathcal{N}(10, 1)$ (informatif d\'ecentr\'e).
\end{exercice}
